% =============================================================================
\section{Amostragem, teorema de Nyquist e aliasing}\label{sec:amostragem}
% =============================================================================

\subsection*{Teoria mínima}

Amostrar é multiplicar por um trem de impulsos
$p(t)=\sum_n\delta(t-nT_s)$. No domínio da frequência isso \textbf{repete} o
espectro a cada $F_s=1/T_s$:
\[
  x_s(t)=x(t)p(t)\ \longleftrightarrow\
  X_s(j\Omega)=\frac{1}{T_s}\sum_{k=-\infty}^{\infty}
  X\bigl(j(\Omega-k\Omega_s)\bigr),\qquad \Omega_s=\frac{2\pi}{T_s}=2\pi F_s .
\]

\begin{teo}[teorema da amostragem]
Se $x(t)$ é limitado em banda a $B$ Hz ($X(j\Omega)=0$ para $|\Omega|>2\pi B$)
e $F_s>2B$, as réplicas \textbf{não se sobrepõem} e $x(t)$ pode ser
reconstruído exatamente por um filtro passa-baixas ideal de corte $F_s/2$. A
frequência $F_N=F_s/2$ é a \textbf{frequência de Nyquist}.
\end{teo}

\paragraph{Aliasing (dobramento).} Se $F_s<2B$, as réplicas se sobrepõem: cada
componente de frequência $f$ aparece na banda digital como
\[
  \boxed{\ f_{\text{aparente}}=\bigl|f-kF_s\bigr|,\quad
  \text{escolhendo }k\text{ tal que }0\le f_{\text{aparente}}\le F_s/2\ }
\]
Mnemônico: \textbf{faixa de Nyquist} $[0,F_s/2]$ com ``espelhamento'' em
múltiplos ímpares de $F_s/2$:
\[
  f_{\text{ap}}=f\ \bmod\ F_s,\qquad
  f_{\text{ap}}\leftarrow F_s-f_{\text{ap}}\ \text{se }f_{\text{ap}}>F_s/2 .
\]
Notas importantes: (i) $f=0$ e $f=F_s$ são indistinguíveis; (ii) $f=F_s/2$
``cai em cima'' da borda da banda; (iii) um tom exatamente em $kF_s$ vira
CC (degenera em constante).

\paragraph{Decimação.} Reduzir a taxa por $M$ (manter 1 de cada $M$ amostras)
leva $F_s\to F_s/M$; para não criar aliasing \textbf{é obrigatório} um filtro
passa-baixas digital (anti-aliasing) com corte $\omega_c=\pi/M$ \emph{antes} de
descartar amostras.

\subsection*{Exercícios resolvidos}

% ---------------------------------------------------------------- 6.1
\exer{Um sinal contém as componentes $200$, $600$, $1300$, $2000$ e $3000$ Hz e
é amostrado a $F_s=4000$ Hz. Qual a frequência aparente de cada componente?}

\textbf{Solução.} $F_s/2=2000$ Hz. Aplicando
$f_{\text{ap}}=|f-kF_s|$:
\[
  \begin{array}{c|c|c}
    f\ \text{[Hz]} & k & f_{\text{ap}}\ \text{[Hz]}\\ \hline
    200  & 0 & 200\\
    600  & 0 & 600\\
    1300 & 0 & 1300\\
    2000 & - & 2000\ (\text{exatamente em Nyquist})\\
    3000 & -1 & |3000-4000|=1000
  \end{array}
\]
Ou seja: $3000$ Hz ``se disfarça'' de $1000$ Hz — é \emph{aliasing}.

\resp{200, 600, 1300, 2000 (borda) e \textbf{1000} Hz (o tom de 3000 Hz
dobrou).}

% ---------------------------------------------------------------- 6.2
\exer{\textbf{(Desafio 01)} O sinal acima é decimado por $M=4$ (nova
$F_s'=1000$ Hz) \emph{sem} filtro anti-aliasing. Determine as frequências
aparentes após a decimação.}

\textbf{Solução.} Agora a banda-base digital é $[0,500]$ Hz:
\[
  200\ \text{Hz}\to200\ \text{Hz};\qquad
  600\ \text{Hz}\to|600-1000|=400\ \text{Hz};\qquad
  1300\ \text{Hz}\to|1300-1000|=300\ \text{Hz}.
\]
Conclusão: as componentes de $600$ e $1300$ Hz sofreram \textbf{aliasing} e
caíram em $400$ e $300$ Hz. Com um FIR anti-aliasing de corte
$\omega_c=\pi/M=\pi/4$ aplicado antes, ambas seriam atenuadas antes da
decimação e o sinal decimado conteria essencialmente apenas a componente de
$200$ Hz.

\resp{$200\to200$, $600\to\textbf{400}$ e $1300\to\textbf{300}$ Hz
(as duas últimas por aliasing).}

% ---------------------------------------------------------------- 6.3
\exer{\textbf{(desafio de reflexão)} Com $F_s=15$ kHz, mostre que um tom de
$9$ kHz aparece em $6$ kHz. E se $F_s=12$ kHz, o que ocorre com o tom de
$6$ kHz?}

\textbf{Solução.} Para $F_s=15$ kHz, $F_s/2=7{,}5$ kHz. Como
$9>7{,}5$, o tom dobra em relação a $F_s$:
\[
  f_{\text{ap}}=|9-15|=6\ \text{kHz}.
\]
Ou seja, um tom de $9$ kHz torna-se \emph{indistinguível} de um de $6$ kHz —
por isso, no Desafio 08, a ``imagem'' da ressonância de $-6$ kHz (que após a
amostragem cai em $+9$ kHz) volta para $6$ kHz e \textbf{soma-se em fase} com
a componente original, inflando o pico do filtro digital.

Para $F_s=12$ kHz: $F_s/2=6$ kHz, então o tom de $6$ kHz cai
\textbf{exatamente em Nyquist} — não há dobramento, mas o tom fica na borda da
banda (na prática: sensível a qualquer imperfeição de filtro e à fase). No
Desafio 08 esse é o caso crítico em que o pico do filtro digital
\emph{dobra} de amplitude.

\resp{$9$ kHz $\to6$ kHz para $F_s=15$ kHz; para $F_s=12$ kHz o tom de $6$ kHz
fica exatamente em $F_s/2$ (Nyquist) — fronteira em que o pico digital dobra.}
